% ============================================================
% meta.tex  --  data diri & judul skripsi
% ============================================================
\usepackage{graphicx}

\newcommand{\judul}{MODEL TEREDUKSI DUA FASA PIPA KALOR NATRIUM UNTUK PENDINGINAN PASIF REAKTOR ANTARIKSA INDONESIA: VALIDASI TERHADAP UJI KRUSTY DAN APLIKASI PADA TERAS 61 PIN 500~kWth}
\newcommand{\judulinggris}{REDUCED-ORDER TWO-PHASE MODEL OF A SODIUM HEAT PIPE FOR PASSIVE COOLING OF THE INDONESIAN SPACE REACTOR: VALIDATION AGAINST THE KRUSTY TEST AND APPLICATION TO A 61-PIN 500~kWth CORE}

\newcommand{\nama}{Muhammad Naufal Tamimi}
\newcommand{\nim}{20/456315/TK/50445}
\newcommand{\tahundaftar}{2020}
\newcommand{\tahun}{2026}

\newcommand{\prodi}{Teknik Nuklir}
\newcommand{\prodien}{Nuclear Engineering}

\newcommand{\kotatanggal}{Yogyakarta, \today}
\newcommand{\tglujian}{tanggal bulan tahun ujian}
\newcommand{\tglpengesahan}{DD BULAN TAHUN}
\newcommand{\tglujianinggris}{Month Date, year}

\newcommand{\pembimbingutama}{Nama Lengkap Pembimbing Utama}
\newcommand{\pembimbingpendamping}{Nama Lengkap Pembimbing Pendamping}
\newcommand{\ketuasidang}{Nama Lengkap Ketua Sidang}
\newcommand{\pengujiutama}{Nama Lengkap Penguji Utama}
\newcommand{\anggotapenguji}{Nama Lengkap Anggota Penguji}

\newcommand{\ketuadepartemen}{Dr. Ir. Alexander Agung, S.T., M.Sc., IPU}
\newcommand{\nipketua}{19720916 199803 1002}

\newcommand{\logougm}{%
  \IfFileExists{fig/Lambang-UGM-hitam.png}%
    {\includegraphics[width=5.5cm]{fig/Lambang-UGM-hitam.png}}%
    {\fbox{\parbox[c][5.5cm][c]{5.5cm}{\centering\small Letakkan logo UGM}}}
}
